\documentclass[a4paper,10pt]{article}
\usepackage[francais]{babel}
\usepackage{url}

\selectlanguage{francais}
\usepackage[utf8]{inputenc}
\usepackage[T1]{fontenc}
\usepackage{babel}
%\usepackage{graphicx}
\usepackage{xspace}
\usepackage{color}
\usepackage{setspace}
\usepackage{geometry}
\usepackage{helvet}
\usepackage{calc}
\usepackage{ifthen}


\pagestyle{empty}

%%%%%
%%%%%

% Package pour la gestion des couleurs
%
\usepackage{color}

% Les commandes pour passer de la version nb à la version couleur 
% (et réciproquement)
%
\newboolean{enNoir}
\setboolean{enNoir}{true}
\newboolean{enCouleur}
\setboolean{enCouleur}{false}
\newcommand{\enNoir}{\definecolor{couleur}{rgb}{0,0,0}
  \setboolean{enCouleur}{false}
  \setboolean{enNoir}{true}}
\newcommand{\enGris}{\definecolor{couleur}{rgb}{0,0,0}
  \setboolean{enCouleur}{false}
  \setboolean{enNoir}{false}}
% Définition du bleu ULP
\newcommand{\enCouleur}{\definecolor{couleur}{rgb}{0,0.2,0.6}
  \setboolean{enCouleur}{true}
  \setboolean{enNoir}{false}}
% La version par défaut en en noir
%
\enNoir

%%%%%%
% Gestion de pdf
%%%%%%
%--------------------------------------
%pour compiler avec latex ou pdflatex
\newif\ifpdf
\ifx\pdfoutput\undefined
  \pdffalse
\else
  \pdfoutput=1
  \pdftrue
\fi
%--------------------------------------
%pour eviter de generer des fonts de type 3
\ifpdf 
\usepackage{pslatex}
\usepackage[pdftex]{graphicx}
%\DeclareGraphicsExtensions{.gif}
\else
\usepackage[dvips]{graphicx}
\fi
%--------------------------------------


\addtolength{\parskip}{\baselineskip}

\newlength{\firstcolumnwidth}
\setlength{\firstcolumnwidth}{1.65in}
\newlength{\inbetween}
\setlength{\inbetween}{5mm}
\newlength{\headermargin}
\setlength{\headermargin}{49mm}

\geometry{left=0.47in+\firstcolumnwidth+\inbetween,right=1in,top=15mm,bottom=20mm,noheadfoot}

\newenvironment{letter}[3][]
{
\thispagestyle{empty}

\newsavebox{\dimargin}
\sbox{\dimargin}{%
  \begin{minipage}[t]{\firstcolumnwidth}
    \vspace{0pt}
    \sffamily
    \begin{flushright}
      %%% Entête du DI
      \newsavebox{\dihead}
      \sbox{\dihead}{%
        \begin{minipage}{\linewidth}
          \flushright
          {\bfseries
            \color{couleur}
            \begin{spacing}{1}
              \large{Laboratoire}
              
              \large{des sciences de l'image,}
            
              \large{de l'informatique et} 

              \large{de la télédétection (LSIIT)} 
        
              \normalsize{UMR 7005} 
            \end{spacing}
            }
          \vspace{-2mm}
          
          \rule{36.5mm}{0.5pt}
          
%          \vspace{-4mm}
          %%% Inclusion du Logo ULP
          \ifthenelse{\boolean{enNoir}}
          {
            \flushright{\includegraphics[height=19mm]{logo-uds-couleur-800x433}}
          }
          {
            \ifthenelse{\boolean{enCouleur}}
            {
              \includegraphics[height=19mm]{}
              }
            {
              \includegraphics[height=19mm]{logo-uds-couleur-800x433}
              }
            }
          %%% Fin inclusion logo
        \end{minipage}%
        }
      \newlength{\diheadheight}
      \newlength{\diheaddepth}
      \newlength{\diheadspace}
      \settoheight{\diheadheight}{\usebox{\dihead}}
      \settodepth{\diheaddepth}{\usebox{\dihead}}
      \setlength{\diheadspace}{\diheadheight+\diheaddepth}
      %
      %
      %%% Coordonnées de l'expéditeur
      \newsavebox{\sender}
      \sbox{\sender}{%
        \begin{minipage}{\linewidth}
          \flushright
          \begin{spacing}{1.5}
            
            \textbf{\small{\fromname}}

            \textbf{\footnotesize{\fonction}}
            
            \scriptsize{Tel. : \phoneespla}
            
            \scriptsize{\phoneillkirch}
            
            \scriptsize{\emailadd}
          \end{spacing}
        \end{minipage}%
        }
      \newlength{\senderheight}
      \newlength{\senderdepth}
      \newlength{\senderspace}
      \settoheight{\senderheight}{\usebox{\sender}}
      \settodepth{\senderdepth}{\usebox{\sender}}
      \setlength{\senderspace}{\senderheight+\senderdepth}
      %
      %
      %%% Coordonnées du DI
      \newsavebox{\difoot}
      \sbox{\difoot}{%
        \begin{minipage}{\linewidth}
          \flushright
          \begin{spacing}{1.25}
            \footnotesize
            
            \textbf{LSIIT-ICPS}
            
            
            {Pôle API -Bd S. Brant}
            
            {F-67400 Illkirch}
            
            {Tél. : +33 (0) 3 90 24 45 42}
            
            {Fax : +33 (0) 3 90 24 45 47}
            
            {http\string://lsiit-cnrs.unistra.fr}
          \end{spacing}
        \end{minipage}%
        }
      \newlength{\difootheight}
      \newlength{\difootdepth}
      \newlength{\difootspace}
      \settoheight{\difootheight}{\usebox{\difoot}}
      \settodepth{\difootdepth}{\usebox{\difoot}}
      \setlength{\difootspace}{\difootheight+\difootdepth}
      %
      %
      \newlength{\devantsender}
      \setlength{\devantsender}{57mm}
%
      \newlength{\devantdifoot}
%
      \setlength{\devantdifoot}{\textheight-\diheadspace-\devantsender-\senderspace-\difootspace-8mm-13pt}
%
      \vspace{8mm}

      \usebox{\dihead}

      \vspace{\devantsender}

      \usebox{\sender}

      \vspace{\devantdifoot}

      \usebox{\difoot}
    \end{flushright}
  \end{minipage}%
}

\begin{spacing}{1}

  \noindent
  \begin{picture}(0,0)
    \setlength{\unitlength}{\firstcolumnwidth+\inbetween}
    \put(-1,0){\usebox{\dimargin}}
  \end{picture}%
%
\begin{minipage}[t]{\headermargin}

\vspace{26mm}

~
\end{minipage}%
\begin{minipage}[t]{\linewidth-\headermargin}
  %Strasbourg, le 30 octobre 2012
  Strasbourg, le \today
\end{minipage}

\noindent
\begin{minipage}[t]{\headermargin}
\vspace{35mm}
~
\end{minipage}%
%-----------------------------------------
% Champ destinataire (#2=To/A , #3 <nom>)   
%-----------------------------------------
\ifthenelse{\equal{#3}{}}{%
\begin{minipage}[t]{\linewidth-\headermargin}

\vspace{2mm}

\end{minipage}
}%
{%
\begin{minipage}[t]{\linewidth-\headermargin} #2
  
  

  \vspace{2mm}

  #3
}
\end{minipage}
\ifthenelse{\equal{#1}{}} {} {Subject : {#1}}
\vspace{20mm}
}
{
%\vspace{20mm}

\newlength{\taillesignature}
\settowidth{\taillesignature}{\fromname}
\newlength{\possignature}
\setlength{\possignature}{0.667\linewidth-0.5\taillesignature}

\hspace{\possignature}\fromname


\end{spacing}
}



\newcommand*{\name}[1]{\def\fromname{#1}}
\name{}
\newcommand*{\fonction}[1]{\def\fonction{#1}}
\newcommand*{\telephoneesplanade}[1]{\def\phoneespla{#1}}
\telephoneesplanade{}
\newcommand*{\telephoneillkirch}[1]{\def\phoneillkirch{#1}}
\telephoneillkirch{}
\newcommand*{\email}[1]{\def\emailadd{#1}}
\email{}
\newcommand*{\objet}[1]{\def\emailadd{#1}}





\selectlanguage{francais}
\begin{document}
\name{Stéphane \textsc{Genaud}}
\fonction{Maître de conférences (HDR)}
\telephoneillkirch{+33 (0)3688 54542}
\email{genaud@unistra.fr}


\begin{letter}[Recommendation de Jonathan Rouzaud-Cornabas]{au}{CoCNRS Section 6}

\textit{Concours CNRS 06/03,  thème  prioritaire  ``calcul haute  performance,
  cloud-computing''}

Madame, Monsieur,


Je connais Jonathan depuis que nous travaillons ensemble sur le projet ANR SONGS
(11-INFRA-13) démarré en  décembre 2012.  Ce projet  poursuit l'objectif général
d'enrichir  l'outil de  simulation  à évènements  discrets \textsc{SimGrid},  en
permettant la  modélisation de  systèmes informatiques  de divers  domaines (les
systèmes pair-à-pair, le  HPC, les clouds, ...). J'anime dans  ce projet le work
package \textit{Clouds}.  L'objectif  est de proposer des  modèles permettant de
prédire le comportement  des applications sur ces  grandes infrastructures. Nous
nous  intéressons principalement  aux  infrastructures dont  on  peut louer  les
ressources  physiques,  appelées  communément  clouds  IaaS.   

Cette entreprise suppose de modéliser  de nombreux mécanismes, tels l'allocation
des  machines virtuelles  sur  les  machines physiques  ou  le comportement  des
systèmes de  stockage (performance  du transfert  des données  et des  images de
machines virtuelles),  mais également  d'inclure les  éléments utilisés  pour la
facturation. Nous  nous appuyons  sur les  modèles éprouvés  de \textsc{SimGrid}
(comme par exemple les modèles de réseau)  ainsi que de son moteur de simulation
comme  brique  de base.  Néanmoins,  la  modélisation  des clouds  nécessite  de
construire une architecture complexe comportant de très nombreux nouveaux
éléments.

Jonathan  a  apporté,   dans  ce  travail,  une  contribution   majeure,  qui  a
impressionné  tout le  groupe.   Il  a réussi  a  développer  en quelques  mois,
l'essentiel du  simulateur que nous  souhaitions avoir. Cette réalisation  lui a
demandé un travail de bibliographie et d'enquêtes minutieuses pour déterminer le
fonctionnement  précis et  les paramètres  réalistes permettant  de simuler  des
infrastructures telles  celles d'Amazon  EC2. Il a  ensuite fallu  concevoir une
structuration logicielle suffisamment générique capable  de décrire le système à
la fois du  côté du client (qui voit  les clouds comme des boîtes  noires) et du
côté  fournisseur   de  l'infrastructure.    La  conception   de  l'architecture
logicielle  est également  ingénieuse en  offrant l'évolutivité  nécessaire pour
s'adapter aux  modes potentiellement différents  de fonctionnement de la multitude
de fournisseurs d'infrastructures, qu'ils soient de type cloud privé ou cloud public. 
Ce développement a été présenté au workshop annuel de l'INRIA-Illinois Joint 
Laboratory en novembre 2012.

La liste de  publications de Jonathan montre par ailleurs  son dynamisme dans la
valorisation  de  ses résultats,  que  ce  soit  sur  son précédent  domaine  de
recherche concernant la  sécurité des systèmes distribués, ou  maintenant sur le
thème de l'ordonnancement des machines  virtuelles dans les clouds, complété par
ce travail sur la simulation des clouds  IaaS. Jonathan est auteur de 3 articles
dans des journeaux internationaux, 7 conférences internationales, et 2 workshops
d'une conférence  majeure du domaine. Dans  le contexte de l'équipe  dynamique à
laquelle il participe, plusieurs autres articles sont déjà en préparation sur le
travail décrit précédemment.

En  conclusion, Jonathan  m'a montré  des capacités  de travail  et de  ténacité
remarquables.  Il  a aussi démontré  qu'il avait été  capable de s'adapter  à un
nouvel  environnement   (de  l'Université   d'Orléans  à l'ENS   Lyon)  très
rapidement, gr\^ace à son investissement  sur des développements logiciels lourds,
qui prennent habituellement plus de  temps pour être valorisés. Cette expérience
réussie en toute  autonomie m'amène à recommander  chaleureusement Jonathan pour
un    poste    de   CR2,    particulièrement    sur    le   thème    prioritaire
\textit{cloud-computing}.





\end{letter}
\begin{flushright}
\includegraphics[width=.30\textwidth]{signgenaud.jpg}
\end{flushright}
%\newpage
%\bibliographystyle{plain}dd%\bibliography{julien_blaise_ED}


\end{document}
